% Options for packages loaded elsewhere
\PassOptionsToPackage{unicode}{hyperref}
\PassOptionsToPackage{hyphens}{url}
%
\documentclass[
]{article}
\usepackage{amsmath,amssymb}
\usepackage{lmodern}
\usepackage{ifxetex,ifluatex}
\ifnum 0\ifxetex 1\fi\ifluatex 1\fi=0 % if pdftex
  \usepackage[T1]{fontenc}
  \usepackage[utf8]{inputenc}
  \usepackage{textcomp} % provide euro and other symbols
\else % if luatex or xetex
  \usepackage{unicode-math}
  \defaultfontfeatures{Scale=MatchLowercase}
  \defaultfontfeatures[\rmfamily]{Ligatures=TeX,Scale=1}
\fi
% Use upquote if available, for straight quotes in verbatim environments
\IfFileExists{upquote.sty}{\usepackage{upquote}}{}
\IfFileExists{microtype.sty}{% use microtype if available
  \usepackage[]{microtype}
  \UseMicrotypeSet[protrusion]{basicmath} % disable protrusion for tt fonts
}{}
\makeatletter
\@ifundefined{KOMAClassName}{% if non-KOMA class
  \IfFileExists{parskip.sty}{%
    \usepackage{parskip}
  }{% else
    \setlength{\parindent}{0pt}
    \setlength{\parskip}{6pt plus 2pt minus 1pt}}
}{% if KOMA class
  \KOMAoptions{parskip=half}}
\makeatother
\usepackage{xcolor}
\IfFileExists{xurl.sty}{\usepackage{xurl}}{} % add URL line breaks if available
\IfFileExists{bookmark.sty}{\usepackage{bookmark}}{\usepackage{hyperref}}
\hypersetup{
  pdftitle={RMarkdown\_FirstAssignment\_Group5},
  pdfauthor={Adrian Enrico Rodriguez, Gian Carlo Fernandez, Iris Salangad, John Kyle Lluz},
  hidelinks,
  pdfcreator={LaTeX via pandoc}}
\urlstyle{same} % disable monospaced font for URLs
\usepackage[margin=1in]{geometry}
\usepackage{graphicx}
\makeatletter
\def\maxwidth{\ifdim\Gin@nat@width>\linewidth\linewidth\else\Gin@nat@width\fi}
\def\maxheight{\ifdim\Gin@nat@height>\textheight\textheight\else\Gin@nat@height\fi}
\makeatother
% Scale images if necessary, so that they will not overflow the page
% margins by default, and it is still possible to overwrite the defaults
% using explicit options in \includegraphics[width, height, ...]{}
\setkeys{Gin}{width=\maxwidth,height=\maxheight,keepaspectratio}
% Set default figure placement to htbp
\makeatletter
\def\fps@figure{htbp}
\makeatother
\setlength{\emergencystretch}{3em} % prevent overfull lines
\providecommand{\tightlist}{%
  \setlength{\itemsep}{0pt}\setlength{\parskip}{0pt}}
\setcounter{secnumdepth}{-\maxdimen} % remove section numbering
\ifluatex
  \usepackage{selnolig}  % disable illegal ligatures
\fi

\title{RMarkdown\_FirstAssignment\_Group5}
\author{Adrian Enrico Rodriguez, Gian Carlo Fernandez, Iris Salangad,
John Kyle Lluz}
\date{7/11/2021}

\begin{document}
\maketitle

\hypertarget{aside-from-the-reasons-mentioned-in-class-why-do-we-need-to-study-and-analyze-data}{%
\subsection{Aside from the reasons mentioned in class, why do we need to
study and analyze
data?}\label{aside-from-the-reasons-mentioned-in-class-why-do-we-need-to-study-and-analyze-data}}

\hypertarget{rico}{%
\subsubsection{Rico}\label{rico}}

Data science is imperative not only in identifying underlying problems
within communities, but they are also predecessors to coming up with
viable solutions to them. This is reflected in the report of the Senate
Economic Planning Office (2017) regarding solid waste management in the
Philippines. Through the meticulous analysis of data dealing with amount
of waste generated, source of wastes among sectors, and composition of
solid waste, the organization was able to identify issues surrounding
solid waste management in the country. As a result, solutions were
proposed to manage waste more efficiently, and effect positive change.

\hypertarget{gian}{%
\subsubsection{Gian}\label{gian}}

For me, data science and analysis has always been part and parcel for a
society to grow and move forward. To make sense of facts and figures and
use it as a logical foundation is the simplest pattern in innovating in
any field. Descriptive Analysis such as what scientists use now to count
COVID-19 cases is what helped the policy-makers to adjudicate the
quarantine limitations on a specific area. Inferential Analysis by the
Philippine Statistics Authority serves as the national basis of data or
census of the citizens. Predictive Analysis by the weather experts is
what warns the population for incoming threats, and Causal Analysis made
by researchers are what innovates our society. Data Science and Analysis
is functional in almost every, if not all fields of our community that
are beneficial for our comfort and well-being. And studying or using
data analysis is helpful, necessary even, to enter any field and become
a functional member of society.

\hypertarget{lluz}{%
\subsubsection{Lluz}\label{lluz}}

Data analysis can help in the industry, where production costs can be
minimized and product quality and production can be maximized to give
both positive results to the producers and consumers. Oftentimes, when
the producers simply want to gain maximum profits but neglect their
customers, the data would show that there is a decrease in product
quality since producers would tend to cut corners. This in turn over
time can lead to a poor company and consumer relationship.

\hypertarget{salangad}{%
\subsubsection{Salangad}\label{salangad}}

These days, studying and analyzing data is already considered a
necessity. With this, we acquire problem solving skills. We tend to
think logically and solve our issues and problems in a proper manner in
a way that is useful and more efficient. Aside from that, we can see
that analytics is everywhere. Whenever we use social media, we can see
the analytics and it is important to know what these analytics are all
about. Knowing analytics is also useful when starting a business to know
the trends and patterns. Today in the 21st century, data and analytics
are used and created every single second. Being knowledgeable about
these data is useful to know the current trends and interests of the
consumers.

\hypertarget{cite-at-least-one-instance-or-application-where-data-science-was-used-and-contributed-to-policy-making-day-to-day-decisions-knowing-market-trends-and-scientific-research.}{%
\subsection{Cite at least one instance or application where data science
was used and contributed to policy-making, day-to-day decisions, knowing
market trends, and scientific
research.}\label{cite-at-least-one-instance-or-application-where-data-science-was-used-and-contributed-to-policy-making-day-to-day-decisions-knowing-market-trends-and-scientific-research.}}

\hypertarget{rico-1}{%
\subsubsection{Rico}\label{rico-1}}

Data science is handy in various sectors and occasions. In policy-making
for example, the aforementioned discussion of the Senate Economic
Planning Office (2017) regarding solid waste management was able to
produce concrete solutions. One of which is the filing of various laws
to further strengthen the Republic Act No.~7160, which focuses on better
controlling waste and allocating resources among local government units
in order to incur minimum waste. With regards to day-to-day decisions,
data science and statistics is a major contributor in choosing which
products to buy in certain markets. Using brands of toothpaste as an
example, consumers would choose products that are more preferred by
dentists than the ordinary brand. The popular brand ``Colgate'' for
example, advertises their toothpaste to be recommended by 9 out of 10
dentists to be better than other brands. Not only does this increase
their credibility, but it also affects consumer preference.
https://ejournals.ph/article.php?id=12731 In terms of market trends,
perhaps one of the most important applications of data science,
statistics can be seen and practiced in the stock market. One of the key
terminologies to be familiar with to find success in stocks is
``technical analysis''.

\hypertarget{gian-1}{%
\subsubsection{Gian}\label{gian-1}}

Data science and statistics, although often overlooked and undervalued,
is probably the biggest aid in governing. Good governance is choosing
the decisions that benefit the society the most with full consideration
towards its members. For our institutions to produce the results that
meet the needs of the society, they should make use of all the available
resources at their disposal, mostly raw data. A prime example of this
would be the City Mayor of Pasig Vico Sotto, who himself gives
importance and pride to a `data-driven policy making governance.'
Getting a hold of the population's situation through a comprehensive
survey may help the local government in making decisions for city-wide
policies, future government-funded projects, and several programs such
as the Social Amelioration Program (SAP).

\hypertarget{lluz-1}{%
\subsubsection{Lluz}\label{lluz-1}}

Outweighing the Pros and Cons of Trophy Hunting

{[}Pro{]}Conservation groups have used this as an advantage to help fund
their operations

{[}Pro{]}Money gained is used as a means to buy guards and defenses for
protecting wildlife against poachers that only gain self-benefit.

{[}Con{]} Funds that are not well managed are often dealt with
corruption

Conclusion: With data analysis, information about these and other
decisions with major effects can be studied and dealt with the best
course of action and not just decisions from a day-to-day basis.
Continuing about the topic. Since Trophy hunting is still seen as an
illegal and inhumane act, further research should be done to gather more
data and create a better conclusion of what are the pros and cons of
doing so. Once the necessary information is analyzed, there is also a
need to take action of what must be done to get the best results based
on the research and thus, there is a need to implement a policy to
ensure that the best course of action is applied as much as possible.
With such policies, there will surely be a change in the trend of the
market in the field of Trophy hunting and thus, those with authority
with regards to the matter shall adapt based on the presented
information of the various data presented.

\hypertarget{salangad-1}{%
\subsubsection{Salangad}\label{salangad-1}}

Scientists from many disciplines are confronted with a lot of
information and are optimistic that the data will help them achieve
their scientific aims. For instance, geneticists have linked genetic
data to vast databases of people's habits and diseases due to the new
sequencing technology that allows for high-resolution genetic sequencing
on a wide scale. These discoveries may aid academics in their research
of the human genome, enabling them to get a greater understanding of how
the genome develops and governs visible characteristics.

\hypertarget{propose-at-least-one-data-science-topic-that-you-want-to-pursue-have-a-broad-description-of-the-topic-describe-the-availability-of-the-data-what-kinds-of-statistical-methods-you-think-you-will-need-and-who-would-benefit-from-this-study.}{%
\subsection{Propose at least one data science topic that you want to
pursue: Have a broad description of the topic, describe the availability
of the data, what kinds of statistical methods you think you will need,
and who would benefit from this
study.}\label{propose-at-least-one-data-science-topic-that-you-want-to-pursue-have-a-broad-description-of-the-topic-describe-the-availability-of-the-data-what-kinds-of-statistical-methods-you-think-you-will-need-and-who-would-benefit-from-this-study.}}

\hypertarget{rico-2}{%
\subsubsection{Rico}\label{rico-2}}

In the Philippines, the primary source of energy is coal - a form of
non-renewable energy that emits several harmful gases in the atmosphere.
With that, there is a call to explore new methods of providing energy
towards Filipino homes with minimal dangers. This is where renewable
energy comes into play. Currently, the country's energy mix is made up
of 47\% coal, 22\% natural gas, 6.2\% oil, and about 24\% renewable
energy. With the given data and others related to it, we may give our
analysis on it and see how much farther we can go with renewable energy.
This requires both descriptive and inferential statistics as it touches
on describing the data and drawing conclusions and generalizations.

\hypertarget{gian-2}{%
\subsubsection{Gian}\label{gian-2}}

A topic I'm interested in studying is the data science of motorcycle
related crimes. I live in the City of Mandaluyong and the
Riding-in-tandem Ordinance is implemented in our city. Several datasets
of motorcycle related crimes in every municipality may be available
online. To know whether the said ordinance truly lowered the crime rates
on the applied cities or was it ineffective. If the policy proved
useful, do the neighboring cities need the same implementation, or does
the ordinance such as this can be improved even more. Inferences such as
these may be helpful for keeping the safety of the citizens, and
analyzing available previous data is the key in figuring everything out.

\hypertarget{lluz-2}{%
\subsubsection{Lluz}\label{lluz-2}}

I would like to study the data between using plastic and shifting to
paper. The topic is to analyse how much carbon dioxide is stored within
these papers and the major effect it leads to the environment. Studies
are being done on how much carbon dioxide is retained or released in the
production of papers. With this, I would like to propose and analyze the
data that will be gathered and draw to a conclusion whether or not the
carbon stored in the tree is released in the process of it becoming
paper, and how much carbon is still retained in the paper product
itself. Conclusion: With data analysis, information about these and
other decisions with major effects can be studied and dealt with the
best course of action.

\hypertarget{salangad-2}{%
\subsubsection{Salangad}\label{salangad-2}}

I would like to pursue machine learning because I'd like to become a
machine learning engineer in the future. Machine learning is a data
analysis technique that automates the construction of analytical models.
It's a field of artificial intelligence based on the concept that
computers can learn from data, recognize patterns, and make choices with
little or no human input. There are machine learning techniques that
create data analytics and predictive analytics training for the data
scientists. This method collects information for statistical analysis.
Machine learning examines and digests enormous quantities of data faster
than a person can, enabling companies to comprehend and adapt in real
time to market and societal developments. Retails can benefit from this
study as machine learning helps them to boost revenue and improve
customer experience data.

\begin{figure}
\centering
\includegraphics{/Users/Lenovo/Documents/R/picture.png}
\caption{Alt text}
\end{figure}

\end{document}
